\begin{lemma}[Дополнение независимого множества]
\(I\subseteq V\) независимо тогда и только тогда, когда \(V\setminus I\) —
вершинное покрытие.
\end{lemma}
\begin{proof}
\(\Rightarrow\). Если ребро \(uv\) не покрыто \(V\setminus I\), то \(u,v\in I\) —
два смежных элемента независимого множества, противоречие.
\(\Leftarrow\). Если \(u,v\in I\) смежны, то ребро \(uv\) не имеет концов в
\(V\setminus I\), т.\,е. не покрыто, противоречие.
\end{proof}
\begin{cor}
\(\alpha(G)+\beta(G)=n\).
\end{cor}
\begin{proof}
Отображение \(I\mapsto V\setminus I\) — биекция между независимыми множествами
и покрытиями, \(|V\setminus I|=n-|I|\). Наибольшему \(I\) соответствует
наименьшее покрытие.
\end{proof}
\begin{graphillustration}{Пример: независимое множество $I=\{a,c\}$ (двойной контур) и его дополнение — вершинное покрытие $\{b,d,e\}$. Каждое ребро имеет конец вне $I$.}{fig:03-cover}
\begin{tikzpicture}[font=\small]
\node[gi selected] (a) at (90:1.2) {$a$};
\node[gi vertex] (b) at (162:1.2) {$b$};
\node[gi selected] (c) at (234:1.2) {$c$};
\node[gi vertex] (d) at (306:1.2) {$d$};
\node[gi vertex] (e) at (18:1.2) {$e$};
\draw[gi edge] (a) -- (b) -- (c) -- (d) -- (e) -- (a);
\node[gi note,anchor=west] at (2.2,0.3) {$C_5$: $\alpha=2$, $\beta=3$,\\$\alpha+\beta=5=n$};
\end{tikzpicture}
\end{graphillustration}


\begin{theorem}[Галлаи, 1959]
Если \(\delta(G)>0\) (нет изолированных вершин), то \(\alpha'(G)+\beta'(G)=n\).
\end{theorem}
\begin{proof}
\(\le\). Пусть \(M\) — максимальное паросочетание, \(|M|=\alpha'\), \(U=V\setminus V(M)\)
— свободные вершины, \(|U|=n-2\alpha'\). Так как \(\delta>0\), для каждой
\(u\in U\) выберем инцидентное ей ребро; пусть \(F\) — множество выбранных рёбер,
\(|F|\le|U|\). Тогда \(M\cup F\) — рёберное покрытие:
\[
\beta'\le|M\cup F|\le|M|+|U|=\alpha'+n-2\alpha'=n-\alpha'.
\]
\(\ge\). Пусть \(L\) — минимальное рёберное покрытие, \(|L|=\beta'\),
\(H=(V,L)\). В \(H\) нет изолированных вершин, поэтому в каждой компоненте
\(H\) есть ребро; выберем по одному ребру из каждой компоненты. Рёбра из разных
компонент не имеют общих концов, так что получилось паросочетание \(N\) графа
\(G\) размера \(\comp(H)\). Отсюда, используя \(e(H)\ge n-\comp(H)\),
\[
\alpha'\ge|N|=\comp(H),\qquad \beta'=|L|=e(H)\ge n-\comp(H)\ge n-\alpha'.\qedhere
\]
\end{proof}
\begin{graphillustration}{Обе половины доказательства Галлаи. Слева: максимальное паросочетание $M$ (толстые) и по ребру (штрих) для каждой свободной вершины; $|M\cup F|=2+2=n-\alpha'$. Справа: минимальное рёберное покрытие — лес звёзд; по ребру из каждой компоненты (толстые) дают паросочетание.}{fig:03-gallai}
\begin{tikzpicture}[font=\small]
\begin{scope}
\node[gi vertex] (b) at (1.2,0.5) {$b$};
\node[gi vertex] (a) at (0,0.5) {$a$};
\node[gi free] (c) at (1.2,1.6) {$c$};
\node[gi free] (u) at (1.2,-0.6) {$u$};
\node[gi vertex] (d) at (2.6,0.5) {$d$};
\node[gi vertex] (e) at (3.8,0.5) {$e$};
\draw[gi match] (a) -- (b); \draw[gi match] (d) -- (e);
\draw[gi edge] (b) -- (d);
\draw[gi alternate] (c) -- (b); \draw[gi alternate] (u) -- (b);
\node[gi note] at (1.9,-1.4) {$M\cup F$, $F=\{cb,ub\}$};
\end{scope}
\begin{scope}[shift={(6.6,0)}]
\node[gi vertex] (s) at (0,0.5) {};
\node[gi vertex] (s1) at (-0.8,1.2) {};
\node[gi vertex] (s2) at (-0.8,-0.2) {};
\node[gi vertex] (s3) at (0.9,1.2) {};
\draw[gi match] (s) -- (s1); \draw[gi edge] (s) -- (s2); \draw[gi edge] (s) -- (s3);
\node[gi vertex] (t) at (2.0,0.5) {};
\node[gi vertex] (t1) at (3.0,0.5) {};
\draw[gi match] (t) -- (t1);
\node[gi vertex] (r) at (4.2,1.0) {};
\node[gi vertex] (r1) at (4.2,0) {};
\node[gi vertex] (r2) at (5.1,0.5) {};
\draw[gi match] (r) -- (r2); \draw[gi edge] (r1) -- (r2);
\node[gi note] at (2.2,-1.4) {$H=(V,L)$: $\comp(H)=3$, $e(H)\ge n-3$};
\end{scope}
\end{tikzpicture}
\end{graphillustration}

\begin{tcolorbox}[colback=cbGrayBack,colframe=cbGray!55]
\textbf{Сводка.} Всегда \(\alpha+\beta=n\), \(\alpha'\le\beta\);
при \(\delta>0\) также \(\alpha'+\beta'=n\) и \(\alpha\le\beta'\)
(каждое ребро покрытия покрывает не более одной вершины независимого множества).
Для двудольных графов неравенства превращаются в равенства (Кёниг).
\end{tcolorbox}
\begin{example}
Для звезды \(K_{1,3}\): \(\alpha=3,\ \beta=1,\ \alpha'=1,\ \beta'=3\); обе суммы равны \(4\).
Для \(C_5\): \(\alpha=2,\ \beta=3,\ \alpha'=2,\ \beta'=3\).
\end{example}
